\FloatBarrier
\section{Moller cone}
\label{sec:moller-cone}

The remaining downstream contacts recorded in this historical study are resolved by the subsequent socket
revision in Section~\ref{sec:downstream-joints}.

The tenth beamline study renames main-cone to moller\_cone in nine source and runtime GXML variations and
nine GEMC2 SQLite CAD rows. The beamline\_W tungsten material, color, passive 180-degree y rotation, and
run-dependent positions remain unchanged. Five placements use cone z = 1589.27 mm and overlap the
FTInnerShield axial interval; four use z = 1273.27 mm and are already separated from the shield. The current
smooth-flange FTInnerShield is unchanged. The raw engineering cone is archived before replacement.

The original contains 15,782 triangles and 302 boundary edges in eight surface loops: five on the outer
tapered wall and three on the inner cylindrical wall. Its open surface has no valid enclosed material
volume. Filling those openings follows the measured outer cone and inner cylinder profiles, preserving every
original face and vertex before clearance work. The closed hole-filled reference has 16,688 triangles. The
large central beam passage remains open. A full nominal axisymmetric reconstruction is avoided because the
engineering outer rings are coarsely tessellated; their existing bounds and facets provide the shape
reference.

The existing entrance counterbore, about 39.98595 mm radius, is widened with a protected 128-sided
cylindrical cut of radius 40.200001 mm. The cut ends at local cone z = -53.53 mm, 0.2 mm beyond the
overlapping shield tip at -53.73 mm. A 0.2 mm nominal radial allowance above the shield radius accounts for
the polygonal opening. Independent clipped-triangle bounds give a minimum cone opening radius of 40.187892
mm and maximum shield radius 40.000001 mm, proving at least 0.187891 mm radial separation throughout the
shared axial interval. Collision checks and bidirectional vertex containment find no intersection in all
nine placements; the nearest sampled vertex-to-surface distance is 0.199938 mm in the overlapping group and
199.109193 mm in the axially separated group.

The clearance reference has 11,616 triangles; protected simplification retains 1,964 joint, shoulder, and
corner facets and produces 5,000 triangles, 68.32\% fewer than the engineering input. File size falls from
789,184 to 250,084 bytes. The final mesh is one closed manifold component of genus one, with no detected
self-intersections or zero-area facets and exactly unchanged outer bounds. Minimum triangle altitude is
0.000380430 mm. Maximum sampled bidirectional displacement from the clearance reference is 0.020048 mm,
below the 0.05 mm shape guard. The original input reproduces the final STL and both reference meshes byte
for byte.

Hole-filled reference volume is 5835559.232219 mm$^3$; final volume is 5825708.437762 mm$^3$. Counterbore
relief and simplification remove 0.168806\% tungsten relative to the closed hole-filled reference.
Simplification alone changes clearance-reference volume by -0.0000002660\%. No original mass difference is
inferred from the defective open input. Fifty-five independent material probes test the filled opening
centres and the retained central beam passage.

\begin{table}[htbp]
\centering
\caption{Original, repaired clearance reference, and final Moller cone.}
\label{tab:moller-mesh}
\begin{tabular}{lrrr}
\toprule
Quantity & Original & Reference & Final \\
\midrule
Triangles & 15,782 & 11,616 & 5,000 \\
Bytes & 789,184 & 580,884 & 250,084 \\
Boundary edges & 302 & 0 & 0 \\
Volume (mm$^3$) & undefined & 5825708.453 & 5825708.438 \\
\bottomrule
\end{tabular}
\end{table}

Two tracking comparisons isolate the cone in air: raw engineering input versus final mesh and repaired
clearance reference versus final mesh. They use the same fixed GEMC3 binary, beamline\_W material, default
placement, and matched target-origin 5 GeV primaries over 1.1--2.6 degrees and full azimuth. FTFP\_BERT runs
without optics or field, with one thread, default production cuts, and no GUI or event output. Each
comparison uses five alternating order repetitions, seeds 12345--12349, and 200 events per electron, proton,
positive-pion, and neutron run. Each has 52 processes including startup and isolated overlap checks. All 104
processes complete normally and all eighty transport runs report 200 events.

Raw/final median transport ratios are 2.200, 1.739, 1.574, 1.552 for electrons, protons, positive pions, and
neutrons respectively. The raw cone emits 26 GeomSolids1001 and 396 GeomNav1002 exception blocks across the
recorded processes. Its stuck-track recovery and invalid solid prevent attributing these ratios solely to
tessellation reduction. The final mesh emits no Geant4 exceptions. The clean repaired-reference/final
comparison gives ratios 1.638, 1.313, 1.402, 1.301 in the same order, with no exceptions in either case.
Ratios above one mean faster tracking. These isolated timings do not establish full CLAS12 speedup or
physics-response equivalence. Median startup simulation time changes from 0.777853 to 0.744276 s; peak
startup RSS changes from 137.08 to 122.14 MiB.

\begin{table}[htbp]
\centering
\caption{Median transport wall seconds per 200 events: raw input versus final cone.}
\label{tab:moller-upgrade}
\begin{tabular}{lrrr}
\toprule
Particle & Original & Final & Original/final \\
\midrule
e- & 55.95160 & 25.43760 & 2.200 \\
proton & 10.28140 & 5.91253 & 1.739 \\
pi+ & 15.50770 & 9.85140 & 1.574 \\
neutron & 11.84350 & 7.63252 & 1.552 \\
\bottomrule
\end{tabular}
\end{table}

\begin{table}[htbp]
\centering
\caption{Median transport wall seconds per 200 events: repaired clearance-reference control.}
\label{tab:moller-reference}
\begin{tabular}{lrrr}
\toprule
Particle & Reference & Final & Reference/final \\
\midrule
e- & 28.23070 & 17.23620 & 1.638 \\
proton & 5.43417 & 4.13941 & 1.313 \\
pi+ & 8.36232 & 5.96464 & 1.402 \\
neutron & 5.71081 & 4.39053 & 1.301 \\
\bottomrule
\end{tabular}
\end{table}

A separate cone/shield assembly test uses 10,000-point Geant4 overlap checking. The original cone gives five
overlap reports and one solid warning; the final assembly has no warnings or overlaps. Full default beamline
construction completes with no missing meshes or solid warnings, retaining four overlap reports at the
unrelated downstream tube/adapter joints. The earlier shield studies retain their historical overlap
measurements; this cone revision resolves that interference. Semantic comparisons pass all fifteen native
geometry and material variations and fourteen CAD variations against GEMC2.

\FloatBarrier
\begin{figure}[htbp]
\centering
\includegraphics[width=0.9\textwidth]{figures/moller_cone_stl_comparison.pdf}
\caption{Original engineering cone (left) and optimized moller\_cone (right), with matched cameras and
bounds. Surface openings are filled and the shield counterbore is widened; the central beam passage and
outer envelope remain.}
\label{fig:moller-cone}
\end{figure}

\begin{figure}[htbp]
\centering
\includegraphics[width=0.95\textwidth]{figures/moller_cone_joint.pdf}
\caption{Actual triangle sections at y = 0 through the original and final cone with the unchanged FT inner
shield. The enlarged counterbore clears the shield tip; the right panel magnifies the former shoulder
interference. The guaranteed radial gap includes polygon faceting.}
\label{fig:moller-joint}
\end{figure}
